% AY2016 力学 解答例
% 体裁・粒度は 参考/ のPDFに揃える（行間を埋め、最終解答は \ans{} で boxed）。
% 解答・数値・問題は本年度過去問から自力で導いた（東大版は体裁の手本のみ）。
\documentclass[11pt,a4paper]{article}
\usepackage{../../../00_common/preamble}

\usetikzlibrary{decorations.pathmorphing}

\title{力学 AY2016 解答例（専門基礎科目）}
\author{}
\date{}

\begin{document}
\maketitle

% ==================================================================
\section*{第1問〔I〕}

天井に固定された質量 $M$，半径 $r$，慣性モーメント $I=\dfrac{1}{2}Mr^2$ の滑車に糸を
巻きつけ，糸の端に質量 $m$ のおもりをつける。おもりは重力で落下し，滑車は回転する。
重力加速度を $g$，糸の張力を $T$，おもりの加速度（落下を正）を $a$ とする。
糸は伸びず滑車上を滑らないので，滑車の角加速度 $\alpha$ とは $a=r\alpha$ で結ばれる。

\subsection*{(1) おもりの並進運動方程式}
おもりには下向きに重力 $mg$，上向きに張力 $T$ がはたらく。落下方向を正にとると
\[
  \ans{\;m a = mg - T\;}
\]

\subsection*{(2) 滑車の回転運動方程式}
滑車には糸の張力 $T$ が半径 $r$ の腕で回転トルク $Tr$ を与える
（吊り糸や軸からの反力は中心 $O$ まわりにモーメントを生まない）。
中心 $O$ まわりの回転の運動方程式は，慣性モーメント $I=\tfrac12 Mr^2$ を用いて
\[
  I\alpha = Tr
  \quad\Longrightarrow\quad
  \frac{1}{2}Mr^2\,\alpha = Tr.
\]
\[
  \ans{\;\dfrac{1}{2}Mr^2\,\alpha = T r,\qquad
  \text{すなわち}\quad \dfrac{1}{2}M r\,\alpha = T\;}
\]

\subsection*{(3) おもりの加速度}
拘束条件 $a=r\alpha$ より $\alpha=a/r$。これを (2) に代入すると
\[
  T=\frac{1}{2}Mr\,\alpha=\frac{1}{2}Mr\cdot\frac{a}{r}=\frac{1}{2}Ma.
\]
これを (1) の $ma=mg-T$ へ代入して
\begin{align*}
  m a &= mg-\frac{1}{2}Ma
  \;\;\Longrightarrow\;\;
  \left(m+\frac{M}{2}\right)a = mg
  \;\;\Longrightarrow\;\;
  a=\frac{mg}{m+\frac{M}{2}}.
\end{align*}
分母を $2$ 倍して整理すると
\[
  \ans{\;a=\dfrac{2m}{2m+M}\,g\;}
\]

\subsection*{(4) 糸の張力}
(3) で得た $T=\dfrac{1}{2}Ma$ に $a=\dfrac{2m}{2m+M}g$ を代入して
\begin{align*}
  T=\frac{1}{2}M\cdot\frac{2m}{2m+M}\,g
   =\frac{Mm}{2m+M}\,g.
\end{align*}
\[
  \ans{\;T=\dfrac{Mm}{2m+M}\,g\;}
\]

\medskip
検算: 滑車が非常に軽い極限 $M\to0$ では $a\to g$（自由落下），$T\to0$ となり，
おもりが抵抗なく落ちる描像と一致。逆に滑車が非常に重い極限 $M\to\infty$ では
$a\to0$，$T\to mg$ となり，おもりがほぼ静止して張力が重力と釣り合う描像と一致。
また $T=mg-ma=m(g-a)>0$ であり，$0<a<g$ と整合する。✓

% ==================================================================
\clearpage
\section*{第2問〔II〕}

静止系 $XY$ に対し，一定角速度 $\omega$ で（$Z$ 軸まわり反時計回りに）回転する
座標系 $X'Y'$ を考える。$t=0$ で両系が一致し，時刻 $t$ で $X'Y'$ 系は $\omega t$ だけ
回転している。点 P の位置を，静止系で $(x,y)$，回転系で $(x',y')$ と表す。

\subsection*{(1) $x$，$y$ の表示}
回転系の基底ベクトルは，静止系から見て
$\bm{e}_{x'}=(\cos\omega t,\ \sin\omega t)$，
$\bm{e}_{y'}=(-\sin\omega t,\ \cos\omega t)$ である。
点 P の位置は $\bm{r}=x'\bm{e}_{x'}+y'\bm{e}_{y'}$ だから，その成分を読みとって
\[
  \ans{\;
  \begin{aligned}
    x &= x'\cos\omega t - y'\sin\omega t,\\
    y &= x'\sin\omega t + y'\cos\omega t
  \end{aligned}\;}
\]
（回転角 $\omega t$ の回転行列による変換である。）

\subsection*{(2) 回転系の軸方向の速度成分}
(1) を時間微分する（$x',y'$ も時間の関数であることに注意）。
\begin{align*}
  \dot{x} &= \dot{x}'\cos\omega t - x'\omega\sin\omega t
             - \dot{y}'\sin\omega t - y'\omega\cos\omega t,\\
  \dot{y} &= \dot{x}'\sin\omega t + x'\omega\cos\omega t
             + \dot{y}'\cos\omega t - y'\omega\sin\omega t.
\end{align*}
速度ベクトル $(\dot{x},\dot{y})$ の，回転系の軸方向（$\bm{e}_{x'},\bm{e}_{y'}$ 方向）
の成分 $v_{x'},v_{y'}$ を内積でとる。
\begin{align*}
  v_{x'} &= \dot{x}\cos\omega t + \dot{y}\sin\omega t\\
         &= \dot{x}'(\cos^2\omega t+\sin^2\omega t)
            - y'\omega(\cos^2\omega t+\sin^2\omega t)
          = \dot{x}' - \omega y',\\
  v_{y'} &= -\dot{x}\sin\omega t + \dot{y}\cos\omega t\\
         &= \dot{y}'(\sin^2\omega t+\cos^2\omega t)
            + x'\omega(\sin^2\omega t+\cos^2\omega t)
          = \dot{y}' + \omega x'.
\end{align*}
（途中の $\dot{x}',\dot{y}',x',y'$ どうしの交差項は $\sin\omega t\cos\omega t$ が
相殺してすべて消える。）
\[
  \ans{\;
  v_{x'} = \dot{x}' - \omega y',\qquad
  v_{y'} = \dot{y}' + \omega x'\;}
\]

\subsection*{(3) 回転系から見た運動方程式}
質点の絶対加速度を回転系の軸方向に射影するため，(2) の $v_{x'},v_{y'}$ を
さらに時間微分する。回転系で測る加速度成分 $a_{x'},a_{y'}$ は，速度ベクトルを
回転系基底に分解した量をもう一度同じ手続きで微分して得られ，結果は
\begin{align*}
  a_{x'} &= \ddot{x}' - 2\omega\dot{y}' - \omega^2 x',\\
  a_{y'} &= \ddot{y}' + 2\omega\dot{x}' - \omega^2 y'
\end{align*}
となる（$\omega$ は一定）。すなわち絶対加速度の回転系成分は，回転系で見た相対加速度
$(\ddot{x}',\ddot{y}')$ に，コリオリ項 $-2\omega\dot{y}',\,+2\omega\dot{x}'$ と
求心項 $-\omega^2 x',\,-\omega^2 y'$ が加わる。
ニュートンの法則 $m(a_{x'},a_{y'})=(F_x',F_y')$ をこれについて解くと
\begin{align*}
  m\ddot{x}' &= F_x' + 2m\omega\dot{y}' + m\omega^2 x',\\
  m\ddot{y}' &= F_y' - 2m\omega\dot{x}' + m\omega^2 y'.
\end{align*}
ベクトル表記（$\bm{\omega}=\omega\bm{e}_z$）では，外力 $\bm{F}'$ に
コリオリ力 $-2m\,\bm{\omega}\times\dot{\bm{r}}'$ と
遠心力 $-m\,\bm{\omega}\times(\bm{\omega}\times\bm{r}')=m\omega^2\bm{r}'$ が加わる。
\[
  \ans{\;
  \begin{aligned}
    m\ddot{x}' &= F_x' + 2m\omega\dot{y}' + m\omega^2 x',\\
    m\ddot{y}' &= F_y' - 2m\omega\dot{x}' + m\omega^2 y'
  \end{aligned}\;}
\]

\medskip
検算: 外力なし $\bm{F}'=\bm{0}$ かつ静止系で等速直線運動する質点を考えると，
静止系では加速度 $0$ であるべきだから，回転系では見かけの力（遠心力＋コリオリ力）
のみで運動する。上式で $\bm{F}'=\bm0$ とおいた右辺がちょうどその見かけの力に等しく，
形式が整合する。また $\omega\to0$ の極限で
$m\ddot{x}'=F_x',\,m\ddot{y}'=F_y'$ と通常のニュートンの運動方程式に帰着する。✓

\end{document}
